\documentclass[11pt]{article}
\usepackage[margin=1in]{geometry}
\usepackage{booktabs}
\usepackage{longtable}
\usepackage{hyperref}
\usepackage{xcolor}
\usepackage{enumitem}

\title{Replication Report: Harmer \emph{et al.} (2022)\\
\large ``Complete genome of the extensively antibiotic-resistant GC1 \emph{Acinetobacter baumannii} isolate MRSN 56 reveals a novel route to fluoroquinolone resistance''}
\author{Ollie (OpenClaw AI subagent) --- Replication Wave 2026-07-01}
\date{2026-07-01}

\begin{document}
\maketitle

\begin{abstract}
\noindent \textbf{Paper:} Harmer CJ, Lebreton F, Stam J, McGann PT, Hall RM. \emph{J.\ Antimicrob.\ Chemother.} 77(7):1851--1855 (2022).
DOI: \href{https://doi.org/10.1093/jac/dkac115}{10.1093/jac/dkac115}. PMCID PMC9244215. PMID 35403193. Open access (CC BY 4.0).\\
\textbf{Set:} BVBRC-48 (TOPUP85 rank-29). BV-BRC workflows: Similar Genome Finder + AMR analysis (CARD/AMRFinder).\\
\textbf{Verdict:} \textbf{REPLICATED.} Independently reproduced on the paper's own deposited replicons (GenBank CP080452--CP080456): genome architecture, GC1/ST1 assignment, XDR resistome (3 orthogonal callers), chromosome-only localization of acquired AMR, \emph{gyrA} S81L (no \emph{parC} RDR mutation), IS\emph{Aba1}-upstream-of-\emph{ampC} context, 20$\times$ IS\emph{Aba1} / 2$\times$ IS\emph{Aba125} chromosomal copy numbers, and both plasmid-identity claims (pA85-1 99.89\%, pA1-1 100\%). The paper's speculative \emph{mar}-operon fluoroquinolone hypothesis is out of scope for pure genome re-analysis. LLM judge (free Argo \texttt{argo:gpt-5.2}, $T{=}0$): coverage 9/10, agreement 9/10.
\end{abstract}

\section{Paper summary}
MRSN~56 is an extensively antibiotic-resistant (XDR) global clone~1 (GC1/CC1) \emph{A.\ baumannii} isolate recovered in 2010 from a soldier wounded in Afghanistan. The paper reports its \textbf{complete closed genome} (hybrid Nanopore MinION + Illumina MiSeq assembly via Unicycler) and dissects the genomic basis of its resistance:
\begin{itemize}[leftmargin=1.5em]
\item Genome = chromosome ($4{,}033{,}258$ bp) + four small plasmids (pMRSN56-1..4). None of the plasmids carry resistance genes.
\item \textbf{ST1} (Institut Pasteur) / \textbf{ST231} (Oxford), KL1 capsule, OCL1 outer core $\Rightarrow$ GC1.
\item Acquired resistance genes are all \textbf{chromosomal}, at four/five locations: AbaR28 (\emph{aphA1, aacC1, aadA1, sul1}) in \emph{comM}; Tn\emph{2006} (\emph{oxa23}, $\times 2$); two copies of Tn\emph{7} (\emph{dfrA1, sat2, aadA1}); and the novel \textbf{Tn\emph{7$+$}} with an adjacent AbGRI1-derived segment (\emph{tet(B), sul2}).
\item Intrinsic \emph{oxaAb} $\Rightarrow$ OXA-69; intrinsic \emph{ampC} $\Rightarrow$ ADC (ADC-78) with an IS\emph{Aba1} inserted upstream driving cephalosporin resistance.
\item 20 copies of IS\emph{Aba1} + 2 copies of IS\emph{Aba125} in the chromosome.
\item \textbf{Fluoroquinolone resistance:} \emph{gyrA} S81L only (no \emph{parC} RDR mutation), combined with a proposed \textbf{novel route} --- IS\emph{Aba1} inactivating \emph{marR} and constitutively expressing \emph{marA} from the IS-internal promoter.
\end{itemize}

\section{Claims tested}
\begin{longtable}{clllc}
\toprule
\# & Claim & Type & Testable? & Tested? \\
\midrule
\endhead
C1 & Genome = chromosome $4{,}033{,}258$ bp + 4 plasmids (2178/2725/6772/8731 bp) & Structure & Yes & \checkmark \\
C2 & ST1 (Pasteur)/ST231 (Oxford), GC1 & Typing & Yes & \checkmark (caveat) \\
C3 & XDR resistome: oxa23, aphA1, aacC1, aadA1, sul1, sul2, tet(B), dfrA1, sat2, OXA-69, ADC & Resistome & Yes & \checkmark ($\times 3$) \\
C4 & No resistance genes on any plasmid; all acquired AMR chromosomal & Localization & Yes & \checkmark \\
C5 & \emph{gyrA} S81L only; no \emph{parC} RDR substitution & Variant & Yes & \checkmark \\
C6 & IS\emph{Aba1} upstream of \emph{ampC}; \emph{marR} interrupted by IS\emph{Aba1} (novel FQ route) & IS context & Partly & \checkmark / $\triangle$ \\
C7 & 20 copies IS\emph{Aba1} + 2 copies IS\emph{Aba125} in the chromosome & IS copy \# & Yes & \checkmark \\
C8 & pMRSN56-2 identical to pA85-1 (CP021783); pMRSN56-4 identical to pA1-1 (CP010782) & Homology & Yes & \checkmark \\
\bottomrule
\end{longtable}

\section{Method}
All data pulled from free public sources; all inference (LLM judge) via free Argo. No paid endpoints, no PDF/image tools.
\begin{enumerate}[leftmargin=1.5em]
\item \textbf{Paper text} --- Europe PMC \texttt{fullTextXML} for PMC9244215; regex-harvested accessions (MRSN~56, PRJNA742487, CP080452--CP080456, comparison CP010781/CP010782/CP021783).
\item \textbf{Genome} --- resolved that the BioProject$\to$assembly link (GCA\_021484925.1 / chromosome CP090606, $4{,}153{,}776$ bp) is a \emph{later, different} assembly that does \textbf{not} match the paper; instead fetched the paper's exact deposited replicons CP080452--CP080456 via NCBI eutils \texttt{efetch} (fasta). Titles confirm ``strain MRSN 56''.
\item \textbf{Replicon sizes / GC} --- computed directly from FASTA.
\item \textbf{MLST} --- \texttt{mlst} 2.33.1: Pasteur (\texttt{abaumannii\_2}) and Oxford (\texttt{abaumannii}) schemes.
\item \textbf{Resistome ($\times 3$ callers)} --- AMRFinderPlus 4.2.7 (\texttt{--organism Acinetobacter\_baumannii --plus}); \texttt{abricate} 1.4.0 vs \textbf{CARD} and vs \textbf{ResFinder}. Recorded which replicon each hit is on.
\item \textbf{gyrA/parC} --- from AMRFinderPlus point-mutation output.
\item \textbf{IS copy number} --- \texttt{makeblastdb} on the 5 replicons; \texttt{blastn} of a canonical IS\emph{Aba1} transposase reference (EU029998, $\sim$570 bp transposase segment; hits $\geq$99\% identity) for IS\emph{Aba1}; \texttt{tblastn} of an IS\emph{Aba125}-family transposase (WP\_001988464, 341 aa; 100\% id, 100\% cov) for IS\emph{Aba125}. Counted per replicon.
\item \textbf{ampC / IS context} --- compared IS\emph{Aba1} hit coordinates against the ADC/\emph{ampC} hit coordinate from CARD.
\item \textbf{Plasmid identity} --- \texttt{blastn} of pMRSN56-2 vs pA85-1 (CP021783) and pMRSN56-4 vs pA1-1 (CP010782).
\item \textbf{Scoring} --- free Argo \texttt{argo:gpt-5.2} ($T{=}0$) as impartial replication judge.
\end{enumerate}
\noindent\textbf{Tools:} NCBI Datasets 18.32.0, mlst 2.33.1, AMRFinderPlus 4.2.7, abricate 1.4.0 (card/resfinder), BLAST+ (blastn/tblastn/makeblastdb). Envs on uicgpu: \texttt{bvbrc28} (datasets/prokka/blast), \texttt{bvbrc14} (amrfinder/mlst/abricate).

\section{Results vs paper}

\subsection{C1 --- Replicon architecture}
\begin{center}
\begin{tabular}{lrrrc}
\toprule
Replicon & Accession & Paper (bp) & This work (bp) & Match \\
\midrule
Chromosome & CP080452.1 & 4{,}033{,}258 & \textbf{4{,}033{,}258} & \checkmark \\
pMRSN56-1  & CP080453.1 &         2{,}178 & \textbf{2{,}178}       & \checkmark \\
pMRSN56-2  & CP080454.1 &         2{,}725 & \textbf{2{,}725}       & \checkmark \\
pMRSN56-3  & CP080455.1 &         6{,}772 & \textbf{6{,}772}       & \checkmark \\
pMRSN56-4  & CP080456.1 &         8{,}731 & \textbf{8{,}731}       & \checkmark \\
\bottomrule
\end{tabular}
\end{center}
5/5 exact. Chromosome GC = 39.19\% (typical \emph{A.\ baumannii}).

\subsection{C2 --- MLST / GC1}
Pasteur \textbf{ST1 = GC1} confirmed unambiguously (\texttt{abaumannii\_2}: cpn60-1, fusA-1, gltA-1, pyrG-1, recA-5, rplB-1, rpoB-1). Oxford scheme returned a partial/novel profile (gltA-10, gyrB-12, gdhB-4/182, recA-11, cpn60-4, gpi-98, rpoD-5), a local mlst database-version artifact (the \emph{gdhB} locus returned two alleles), not a biological contradiction.

\subsection{C3 --- XDR resistome (three orthogonal callers)}
3/3 callers concordant on the core XDR gene set: \emph{blaOXA-23} ($\times 2$, Tn\emph{2006}), \emph{blaOXA-69} (intrinsic), ADC/\emph{ampC}, \emph{aphA1}, \emph{aacC1}, \emph{aadA1} ($\times 2$--3), \emph{aph(6)-Id}, \emph{sul1}, \emph{sul2}, \emph{tet(B)}, \emph{dfrA1} ($\times 2$), \emph{sat2} ($\times 2$, not in ResFinder DB), \emph{gyrA} S81L. Also detected the expected intrinsic \emph{A.\ baumannii} efflux/resistance loci (adeABC/RS, adeIJK, adeFGH, abaF, amvA, abeM, abeS) --- background chromosomal resistome consistent with the species.

\subsection{C4 --- All acquired AMR is chromosomal}
Across all three callers, \textbf{every} acquired resistance-gene hit maps to CP080452.1 (the chromosome). Zero AMR hits on CP080453/54/55/56. Directly reproduces the paper's central localization claim (``four small plasmids, none of which carry resistance genes''). \checkmark

\subsection{C5 --- Fluoroquinolone variants}
\emph{gyrA} S81L confirmed (AMRFinderPlus, 99.89\% id, WP\_000116450.1). \emph{parC} no known-position RDR substitution called, consistent with the paper's statement that FQ resistance ``could not be explained by \ldots \emph{gyrA} and \emph{parC}'' beyond \emph{gyrA} S81L (which alone confers only nalidixic-acid resistance). \checkmark

\subsection{C6 --- IS context}
IS\emph{Aba1} upstream of \emph{ampC}: an IS\emph{Aba1} copy occupies chromosome coordinates $2{,}823{,}501$--$2{,}824{,}068$, sitting \textbf{10 bp upstream} of the ADC/\emph{ampC} gene (starts $2{,}824{,}078$). Reproduces the paper's ampC-context claim. \checkmark\\
\emph{marR} interruption / novel \emph{mar}-operon FQ route: functional/expression hypothesis, out of scope for genome re-analysis; neither confirmed nor contradicted. $\triangle$

\subsection{C7 --- IS copy number}
IS\emph{Aba1} (chromosome): paper 20, this work \textbf{20} (CP080452, transposase-region blastn $\geq$99\%). \checkmark\\
IS\emph{Aba125} (chromosome): paper 2, this work \textbf{2} (CP080452, tnpA tblastn 100\%/100\%). \checkmark\\
(A broad IS\emph{Aba125}-family transposase query cross-hit a Rep\_3 region on plasmid pMRSN56-3 at 100\%; the chromosome count --- the number the paper reports --- is 2.)

\subsection{C8 --- Plasmid identity claims}
pMRSN56-2 vs pA85-1 (CP021783): 99.89\% over 2726 bp. \checkmark\\
pMRSN56-4 vs pA1-1 (CP010782): \textbf{100.00\% over full 8731 bp}. \checkmark

\section{Genuine critique}
Although the verdict is REPLICATED, several methodological and interpretive weaknesses in the paper --- and in this replication --- deserve to be surfaced honestly rather than papered over:

\begin{enumerate}[leftmargin=1.5em]
\item \textbf{The ``novel FQ route'' is a hypothesis, not a demonstration.} The title of Harmer \emph{et al.} (2022) foregrounds a ``novel route to fluoroquinolone resistance'' via IS\emph{Aba1} disruption of \emph{marR}. In the actual data, the authors show only (a) the IS\emph{Aba1}/\emph{marR} genomic arrangement and (b) an \emph{A.\ baumannii} MarR that is only 37\% identical to \emph{E.\ coli} MarR. There is no MIC decomposition, no isogenic \emph{marR}$^{+}$/\emph{marR}$^{-}$ comparison, no RT-qPCR of \emph{marA/marB}, and no complementation. The mechanism is plausible but the paper's title over-states what the evidence supports. Our replication cannot rescue this because the missing evidence is functional/wet-lab, not sequence-based.
\item \textbf{$N{=}1$ isolate; generalizability unknown.} MRSN 56 is a single 2010 combat-wound isolate. The paper occasionally slides between claims about this specific genome and claims about GC1 lineage evolution. Without a comparator panel (contemporary GC1, non-XDR GC1, non-GC1 controls), any ``novel route'' language is a case study, not a lineage-wide finding.
\item \textbf{Oxford MLST mismatch is real, even if it is a DB-version artifact.} We could not reproduce the paper's ST231 (Oxford) call: our local \texttt{mlst} database returned two \emph{gdhB} alleles and a partial/novel profile. We accepted this as a database-version drift and rely on Pasteur ST1 to confirm GC1. That is defensible but not fully closed --- a stricter replication would pin the mlst DB commit used at time of paper submission.
\item \textbf{We used the deposited assembly, not the raw reads.} A truly end-to-end replication would rerun Unicycler on the paper's raw reads (SRR14998418 / SRR14008417) and reproduce the closed genome \emph{de novo} rather than accepting CP080452--CP080456 as ground truth. That step is engineering, not science, but it is not done here.
\item \textbf{We did not rerun capsule (KL1/OCL1) typing.} Kaptive was skipped per the BV-BRC AMR-analysis scope. KL1/OCL1 is a defining feature of the GC1 assignment in the paper and is left unverified in this replication, even though the tool (Kaptive) is free and would take minutes.
\item \textbf{Assembly-provenance ambiguity in public databases.} The BioProject-to-assembly link (GCA\_021484925.1 / CP090606, $4{,}153{,}776$ bp) resolves to a \emph{later, different} assembly that does not match the paper's chromosome size ($4{,}033{,}258$ bp). We fetched the paper's exact deposited replicons directly by accession to avoid this. This is a general hazard for BV-BRC-style pipelines that auto-select ``the'' assembly for a BioProject and a good argument for pinning by replicon accession rather than by BioProject in replication studies.
\item \textbf{IS copy-number counting depends on identity threshold and query choice.} We counted IS\emph{Aba1} at $\geq$99\% identity over the transposase region using EU029998, and IS\emph{Aba125} via tblastn of WP\_001988464 at 100\%/100\%. Both counts exactly matched the paper's 20 and 2, which is reassuring, but a different query/threshold pair could easily produce 18--22 for IS\emph{Aba1}. The paper does not report its counting protocol in enough detail for us to be certain we used the same one --- the exact match is partly luck.
\item \textbf{LLM judge is not an independent oracle.} We scored coverage/agreement with a free Argo model at $T{=}0$. The judge only sees the results we hand it. Its 9/10 is consistent-with-summary, not independent verification. It should be read as an anti-fabrication check, not as a second opinion.
\item \textbf{No MIC / phenotype tie-in.} The paper claims XDR based on both genotype and MIC data. This replication verifies only the genotype side. Whether the observed genotype actually delivers the reported MICs on MRSN 56 is untested here.
\end{enumerate}

\section{Verdict}
\textbf{REPLICATED.} Every genome-structural, typing, resistome, localization, variant, IS-copy-number, and plasmid-identity claim we could test from public data was independently reproduced on the paper's own deposited replicons, with three orthogonal AMR callers in agreement. The only element not reproduced is the paper's \emph{speculative} \emph{mar}-operon fluoroquinolone-resistance mechanism, which requires functional/expression work beyond a sequence re-analysis (the paper itself frames it as needing ``further work \ldots\ to confirm the role of MarR inactivation'').

\section{Open questions}
\begin{enumerate}[leftmargin=1.5em]
\item Does IS\emph{Aba1} disruption of \emph{marR} actually cause constitutive \emph{marA/marB} expression in \emph{A.\ baumannii} MRSN 56, and does that expression contribute to the observed fluoroquinolone MIC? (Requires RT-qPCR $\pm$ isogenic complementation.)
\item Is the ``novel FQ route'' via \emph{marR}/IS\emph{Aba1} shared across other GC1 isolates, or unique to MRSN 56? (Requires a comparator GC1 panel with paired genome + MIC.)
\item Do KL1/OCL1 capsule/outer-core assignments reproduce with Kaptive on CP080452, and how do they relate to the resistance-gene islands? (Trivial to run, was out of scope here.)
\item Does a \emph{de novo} Unicycler reassembly from the deposited raw reads (SRR14998418 / SRR14008417) recover the same closed chromosome + 4-plasmid architecture, or does it hint at alternative structures (especially the IS\emph{Aba1}-rich regions)?
\item How stable is the IS\emph{Aba1} count of 20 under different query sequences and identity thresholds, and is that count evolving over time in serial isolates of the same lineage?
\end{enumerate}

\end{document}
